% ============================================================
% preamble.tex — pacchetti e configurazione del documento
% ============================================================

% --- Codifica e lingua ---
\usepackage[utf8]{inputenc}
\usepackage[T1]{fontenc}
\usepackage[italian]{babel}
\usepackage{csquotes}

% --- Layout di pagina ---
\usepackage[a4paper,margin=3cm]{geometry}
\usepackage{setspace}
\onehalfspacing
\interfootnotelinepenalty=10000 % vieta lo spezzamento di una nota a piè di pagina tra due pagine

% --- Matematica ---
\usepackage{amsmath}
\usepackage{amssymb}
\usepackage{amsthm}
\usepackage{mathtools}

% --- Figure e tabelle ---
\usepackage{graphicx}
\usepackage{booktabs}
\usepackage{siunitx}
\usepackage{caption}
\usepackage{subcaption}
\usepackage{tikz}
\usetikzlibrary{positioning}

% --- Algoritmi (utile per Cap. 2-3-6: eigensolver, message passing, training loop) ---
\usepackage{algorithm}
\usepackage{algpseudocode}
% babel non traduce il float degli algoritmi né le parole chiave di algpseudocode
\floatname{algorithm}{Algoritmo}
\numberwithin{algorithm}{chapter}
\algrenewcommand\algorithmicrequire{\textbf{Input:}}
\algrenewcommand\algorithmicensure{\textbf{Output:}}
\algrenewcommand\algorithmicfor{\textbf{per}}
\algrenewcommand\algorithmicdo{\textbf{esegui}}
\algrenewcommand\algorithmicend{\textbf{fine}}
\algrenewcommand\algorithmicreturn{\textbf{restituisci}}

% --- Bibliografia ---
\usepackage[
    backend=biber,
    style=numeric,
    sorting=nyt
]{biblatex}
\addbibresource{bibliography/references.bib}

% --- Riferimenti incrociati e link (va caricato quasi per ultimo) ---
\usepackage{hyperref}
\hypersetup{
    colorlinks=true,
    linkcolor=black,
    citecolor=black,
    urlcolor=blue,
    bookmarksnumbered=true
}
\usepackage[capitalize,noabbrev]{cleveref}
\crefname{chapter}{Capitolo}{Capitoli}
\Crefname{chapter}{Capitolo}{Capitoli}
\crefname{section}{Sezione}{Sezioni}
\Crefname{section}{Sezione}{Sezioni}
\crefname{algorithm}{Algoritmo}{Algoritmi}
\Crefname{algorithm}{Algoritmo}{Algoritmi}
\crefname{figure}{Figura}{Figure}
\Crefname{figure}{Figura}{Figure}
\crefname{table}{Tabella}{Tabelle}
\Crefname{table}{Tabella}{Tabelle}
\crefname{equation}{Equazione}{Equazioni}
\Crefname{equation}{Equazione}{Equazioni}

% --- Ambienti per definizioni/teoremi/proposizioni (utili nei Cap. 2-3) ---
% Gli ambienti condividono la numerazione della sezione "definition"; senza
% aliascnt, cleveref userebbe il nome del contatore e stamperebbe "Definition"
% anche per proposizioni, teoremi, esempi e osservazioni.
\usepackage{aliascnt}
\theoremstyle{definition}
\newtheorem{definition}{Definizione}[chapter]
\newaliascnt{proposition}{definition}
\newtheorem{proposition}[proposition]{Proposizione}
\aliascntresetthe{proposition}
\newaliascnt{theorem}{definition}
\newtheorem{theorem}[theorem]{Teorema}
\aliascntresetthe{theorem}
\newaliascnt{example}{definition}
\newtheorem{example}[example]{Esempio}
\aliascntresetthe{example}
\newaliascnt{remark}{definition}
\newtheorem{remark}[remark]{Osservazione}
\aliascntresetthe{remark}
\renewcommand{\proofname}{Dimostrazione}
\crefname{definition}{Definizione}{Definizioni}
\Crefname{definition}{Definizione}{Definizioni}
\crefname{proposition}{Proposizione}{Proposizioni}
\Crefname{proposition}{Proposizione}{Proposizioni}
\crefname{theorem}{Teorema}{Teoremi}
\Crefname{theorem}{Teorema}{Teoremi}
\crefname{example}{Esempio}{Esempi}
\Crefname{example}{Esempio}{Esempi}
\crefname{remark}{Osservazione}{Osservazioni}
\Crefname{remark}{Osservazione}{Osservazioni}
